\subsection{Routes}
\label{sec:3_4_ingestion_routes}

\par
The routes ingestion step was the most involved because the routes dataset
stores airport IATA codes, while the assignment requires the country name of both
the source and destination airports. For that reason, I first loaded the routes
file as a CSV DataFrame and then built an airport lookup table from the airports
dataset to translate airport codes into country names (\cref{fig:3-4-routes-csv-read}).

\begin{figure}[H]
	\centering
	\includegraphics[width=0.9\textwidth]{3_4_routes_csv_read.png}
	\caption{Initial CSV load of \texttt{routes.dat}}
	\label{fig:3-4-routes-csv-read}
\end{figure}

\par
My first attempt tried to enrich routes directly by joining the source airport
code against the countries dataset. That produced \texttt{NULL} values because
the join keys were incompatible: route airport codes such as \texttt{AER} and
\texttt{ASF} do not match country ISO codes (\cref{fig:3-4-routes-countries-join-attempt}).

\begin{figure}[H]
	\centering
	\includegraphics[width=0.9\textwidth]{3_4_routes_countries_join_attempt.png}
	\caption{Incorrect direct join attempt between routes and countries,
		resulting in null matches}
	\label{fig:3-4-routes-countries-join-attempt}
\end{figure}

\par
To fix that, I used the airports DataFrame as an intermediate mapping,
selecting the airport IATA code and its country. I then joined the routes
DataFrame against that mapping using \texttt{source\_airport} as the key and
renamed the resulting country column to \texttt{source\_country}. This produced
valid source-country values but was a more involved process than I thought it
would be (\cref{fig:3-4-routes-countries-join-fixed}).

\begin{figure}[H]
	\centering
	\includegraphics[width=0.9\textwidth]{3_4_routes_countries_join_fixed.png}
	\caption{Correct source-country enrichment using the airport-to-country
		mapping}
	\label{fig:3-4-routes-countries-join-fixed}
\end{figure}

\par
I basically repeated the same process for the destination airport and obtained
both \texttt{source\_country} and \texttt{destination\_country}. At that point,
the routes DataFrame contained the required geographic enrichment for each
route (\cref{fig:3-4-routes-source-destination-country}).

\begin{figure}[H]
	\centering
	\includegraphics[width=0.9\textwidth]{3_4_routes_source_destination_country.png}
	\caption{Final enriched routes DataFrame with source and destination country
		columns}
	\label{fig:3-4-routes-source-destination-country}
\end{figure}

\par
The assignment also required removing routes operated by inactive airlines.
I'd already filtered \texttt{airlinesDf} down to active airlines back in
\cref{sec:3_2_ingestion_airlines}, so here I reused those airline IDs to filter out routes operated
by inactive airlines. The raw routes data itself isn't filtered, so this join
was the only thing removing them. This ensured that the final routes dataset
contained only valid airline-operated routes (\cref{fig:3-4-routes-inactive-airlines-filtered}).

\begin{figure}[H]
	\centering
	\includegraphics[width=0.9\textwidth]{3_4_routes_inactive_airlines_filtered.png}
	\caption{Routes after filtering out inactive airlines}
	\label{fig:3-4-routes-inactive-airlines-filtered}
\end{figure}

\par
Finally, I wrote the resulting DataFrame to MongoDB in the
\texttt{practice.routes} collection and read it back to verify the content. The
MongoDB output confirmed that the enrichment columns and the route fields were
stored correctly (\cref{fig:3-4-routes-mongodb-write-read}).

\begin{figure}[H]
	\centering
	\includegraphics[width=0.9\textwidth]{3_4_routes_mongodb_write_read.png}
	\caption{Routes stored in MongoDB and read back for verification}
	\label{fig:3-4-routes-mongodb-write-read}
\end{figure}
